\chapter{Quick start: a certified release in ten minutes}\label{ch:quickstart}
\chapterintro{This chapter takes the motor-mount brief b06 from the task suite to a certified, exported design of
73.7 million elements. You start two processes, read the specification against the brief, climb the resolution
ladder, read the certificate and export the STL with its report. No language model is involved.}

\begin{notebox}[What you need]
A finished installation (Chapter~\ref{ch:install}) with the solver connected, and a GPU with about 24~GB of memory
that no other CUDA program is using. On a smaller GPU, start the backend with a lower ladder cap
(Section~\ref{sec:qs-small}); everything else in this chapter stays the same.
\end{notebox}

\section{The brief}
The brief is a plain-text file, \code{briefs/b06\_motor\_mount.md}:

\begin{quote}\small\sffamily
Plate envelope 160 x 60 x 120 mm (x, y, z), y up, aluminium 6061. It bolts to the machine bed through four M10 bolts on
its bottom face (y = 0) at x = 20 and 140 mm, z = 20 and 100 mm; each bolt clamps through a 24 mm washer. Keep a solid
boss of 24 mm diameter and 10 mm height around each bolt.

The motor sits on a 40 x 40 mm, 5 mm thick solid pad on the top face, centred at x = 80, z = 60. Its mounting pad
carries 2 kN downward and 0.8 kN sideways (+x) from belt tension.

Volume fraction 0.2, 1 mm elements.
\end{quote}

\section{Start the workspace}
\begin{steps}
\item Start the backend and its GPU worker. It listens on \code{127.0.0.1:8765}.
\begin{shell}
python -B scripts/serve.py
\end{shell}
\item In a second terminal, start the web front end:
\begin{shell}
node web/node_modules/vite/bin/vite.js web     # or: cd web && npm run dev
\end{shell}
\item Open \url{http://127.0.0.1:5173} in Chrome. The \gui{GPU} pill at the bottom right reads \emph{starting} until the
worker has loaded CUDA and the solver, then \emph{ready}.
\end{steps}

When the backend session is still empty, the workspace opens the suite task \code{b06\_motor\_mount} with its reference
specification, and shows the brief in the assistant panel. If the session already holds a design (the problem picker
then reads \gui{live session}), choose \gui{b06\_motor\_mount} in the problem picker at the top left.

\section{Read the specification against the brief}
Before solving anything, check that the specification says what the brief says. This is the intent check of
Section~\ref{sec:threechecks}, and no software can do it for you. The \gui{Design intent} tree on the left lists every
item; click one to see its exact values in the \gui{Inspector} and to highlight it in the viewport.
Table~\ref{tab:b06} maps the brief to the specification.

\begin{table}[h]
\caption{Brief b06 and its specification. Coordinates are absolute, in millimetres, with the origin at the domain's
minimum corner.}\label{tab:b06}
\begin{fixtable}{@{}>{\raggedright\arraybackslash}p{0.36\linewidth}>{\raggedright\arraybackslash}X@{}}
Brief & Specification \\ \midrule
envelope 160 x 60 x 120 mm, aluminium 6061 & domain \code{[160, 60, 120]}; material $E = 69000$~N/mm$^2$,
$\nu = 0.33$ \\
four M10 bolts through 24 mm washers on the bottom face & supports \code{bolt\_1}--\code{bolt\_4}: elliptical
\code{patch} regions of size 24 x 24 on face \code{y-}, centred at $(x,z)$ = (20, 20), (20, 100), (140, 20), (140, 100);
all directions fixed \\
solid boss, 24 mm diameter, 10 mm high & keep-in \code{boss\_1}--\code{boss\_4}: cylinders along $y$, radius 12,
span 0--10 \\
40 x 40 mm, 5 mm solid pad on top, centred at x = 80, z = 60 & keep-in \code{pad}: box from (60, 55, 40) to
(100, 60, 80) \\
2 kN down and 0.8 kN in +x on the pad & load \code{motor}: force \code{[800, -2000, 0]}~N on a 40 x 40 patch of face
\code{y+} \\
volume fraction 0.2, 1 mm elements & \code{volume\_fraction} 0.2 (of the design volume); resolution $h = 1$ \\
\end{fixtable}
\end{table}

The \gui{Checks} panel should read \emph{All rules pass}, and the top bar \emph{valid spec}. The ladder at the bottom
then shows four buttons.

\section{Climb the ladder}
The ladder for this specification has four named rungs. The spec's 1~mm elements are the base mesh; the release rung is
the finest grid that fits on the GPU, here 0.25~mm:

\begin{center}\small\sffamily
\begin{tabular}{@{}lrrl@{}}
\toprule
Rung & Element edge & Grid & Elements \\ \midrule
Sketch & 2 mm & 80 x 30 x 60 & 144,000 \\
Preview & 1 mm & 160 x 60 x 120 & 1.15 million \\
Refine & 0.5 mm & 320 x 120 x 240 & 9.2 million \\
Release & 0.25 mm & 640 x 240 x 480 & 73.7 million \\
\bottomrule
\end{tabular}
\end{center}

\begin{steps}
\item Click \gui{Sketch}. Density frames stream into the viewport while it runs. The first sketch of a new
specification builds its multigrid hierarchy and takes a few seconds; on the paper's workstation the cold b06 sketch
took 3.4~s.
\item Click \gui{Preview}, then \gui{Refine}. Each starts from the design of the rung below (marked \gui{$\uparrow$}
on the button) and takes seconds. Figure~\ref{fig:b06refine} shows the refine design.
\item When Refine has converged, click \gui{Optimize release}. It re-opens the optimization on the 73.7-million-element
grid for about 60 iterations. On the paper's workstation this policy used 61\% of the compute of a cold release-rung run
(423~s) and ended within 0.03\% of its compliance \claim{C3.1}. Alternatively, \gui{Release} certifies the upsampled
refine design in about 12 iterations (18~s of compute for b06 in the latency study \claim{C6.4}), but it barely changes
the design (Section~\ref{sec:policies}).
\end{steps}

\begin{figure}[t]
\wideshot{b06_refine.png}
\caption{The motor mount after Refine (9.2 million elements). The result card at the top right shows the checks: five
pass, but \emph{spill monitored} fails because spill monitoring runs only on the release rung, so the card says
\emph{not certified}. That is expected for every rung below Release. Screenshot from \code{scripts/gui\_screenshots.py}.}
\label{fig:b06refine}
\end{figure}

\section{Read the certificate}
When the release finishes, the result card shows the latest design and its certificate. A release is finished only
when the card reads \gui{certified}: all six checks carry a \pass, and the grid is the release grid (here
640 x 240 x 480). Chapter~\ref{ch:ladder} explains each check. The certified design is shown as its $\rho = 0.5$
surface; \gui{Voxels} switches to the raw density field.

\begin{figure}[t]
\centering
\shot[0.47\linewidth]{11_result_card.png}\hfill
\shot[0.47\linewidth]{13_export.png}
\caption{Left: the result card of a certified release, with all six checks passing. Right: the card after export,
with links to the STL and the two report files. These two screenshots come from the bracket of
Chapter~\ref{ch:spec} on a reduced ladder; the card for b06 has the same layout.}
\label{fig:qs-card}
\end{figure}

\section{Export}
Click \gui{Export STL + report}. Three files are written to the session's \code{exports/} folder, under
\code{sessions/}\meta{date-time}\code{/exports/}, and linked from the card (Figure~\ref{fig:qs-card}):

\begin{itemize}
  \item \code{b06\_motor\_mount\_}\meta{job}\code{.stl}: a closed binary STL of the full-resolution $\rho = 0.5$ surface
  (in the latency study the export at release took 0.3--2.2~s \claim{C6.4});
  \item \code{b06\_motor\_mount\_}\meta{job}\code{\_report.md} and \code{\_report.json}: the specification and its hash,
  the rung, grid and solver, iterations and compute time, the compliance, every certificate check, the STL's hash, and
  the hash of the session log at the moment of export (Appendix~\ref{app:certificate}).
\end{itemize}

A design that is not a certified release is exported with \code{\_draft} in its file name, so a draft cannot be
mistaken for a release later.

\begin{tipbox}[You are done]
You now hold a certified design, its STL, a report that ties it to the exact specification and solver settings, and a
session log in \code{sessions/}\meta{date-time}\code{/session.jsonl} from which the whole session can be replayed
(Chapter~\ref{ch:provenance}).
\end{tipbox}

\section{On a smaller GPU}\label{sec:qs-small}
Cap the ladder when you start the backend:
\begin{shell}
python -B scripts/serve.py --ladder-max 1.2e7
\end{shell}
The cap can lower the default of \LadderMax\ elements but never raise it. With this cap b06 gets the rungs 18,000 /
144,000 / 1.15 million / 9.2 million, and the 9.2-million-element rung becomes its release rung. The certificate means
the same on any rung that is the release rung of its ladder: the result is certified on the grid it reports, and the
report states that grid.

\section{The same task with an agent}
The paper ran this task without a person: a Codex agent connected over MCP built the specification from the brief,
asked four clarification questions (answered by a scripted oracle holding the reference reading), climbed the ladder and
exported a certified 73.7-million-element release in 176~s from its first tool call. Its specification equalled the
reference reading in every field, and the log replayed exactly. Chapter~\ref{ch:assistant} shows how to do the same with
the built-in assistant, and Chapter~\ref{ch:mcp} how to connect an agent of your own.
